% !PRINTED 473
\weberpage{473}
%<<CONT
\weq{26}{\frac{c+\partial\omega}{a} = \frac{2\gamma+\delta\omega}{2\alpha+\beta\omega}, \quad \delta \equiv 0 \pmod{2}}

and in addition

\weq{27}{\gamma^2 - 1 + \gamma(2\alpha-\delta) \equiv 0 \pmod{16};}

from (26) and (24) it follows:
\[
\begin{aligned}
2\alpha &= B+y, & 2n\gamma &\equiv -C\\
\beta &= A, & 2n\delta &\equiv -(B-y)
\end{aligned}
\quad (\mathrm{mod}\ 16),
\]
and from this one concludes as above:

9. \gesperrt{If $m \equiv 2 \pmod{4}$, then $\sqrt{2}\,f_1(\omega)^6$ is a class invariant.}

Again all the examples indicate that, if
$m \equiv 4 \pmod{8}$, $\sqrt{2}\,f_1(\omega)^{12}$ is a \gesperrt{class invariant}. But
here too the proof is still lacking. If $m$ is divisible by 8,
no reduction occurs.

In cases 4. to 8. there follows, by 3., a further
reduction to the third root, if $m$ is not divisible by 3.
Accordingly we obtain the following cases:
\[
m \equiv \pm 1 \pmod{3},
\]
\[
\begin{array}{lll}
 & & \text{class invariant}\\
1) & m \equiv 3 \pmod{4} & f(\omega)^2,\\
2) & m \equiv 5 \pmod{8} & f(\omega)^4,\\
3) & m \equiv 1 \ \text{\textquotedbl} & \sqrt{2}\,f(\omega)^2,\\
4) & m \equiv 7 \ \text{\textquotedbl} & \sqrt{2}\,f(\omega),\\
5) & m \equiv 3 \ \text{\textquotedbl} & f(\omega),\\
6) & m \equiv 2 \ \text{\textquotedbl} & \sqrt{2}\,f_1(\omega)^2,\\
7) & m \equiv 4 \ \text{\textquotedbl} & \sqrt{2}\,f_1(\omega)^4,
\end{array}
\]
of which, to be sure, the cases 5) and 7) are inferred only by
 induction.

The class invariants $f(\omega)$ are particularly well suited to
numerical calculation, because among them all they yield the
simplest results. We therefore give first a larger series
of examples, which can easily be derived on the basis of our developments so far and which at the same time
teach the methods that can also be used for further calculations of this kind.

In these calculations one is frequently faced with the task of
decomposing reducible integral rational functions into factors.
Such a decomposition, once found, is of course immediately
%CONT>>
